\documentclass[12pt,a4paper]{article}
\usepackage{amsmath}
\usepackage{graphicx}
\usepackage{color}

%%%% Special package for Latex2html
\usepackage{html}
\usepackage{heqn}

%%%%  Margin modifications --> Geometry package
% \begin{latexonly}
  \usepackage[verbose,hmargin={2cm,2cm},vmargin={2cm,2cm}]{geometry}
% \end{latexonly}

\setcounter{secnumdepth}{1}
\setcounter{tocdepth}{1}

\definecolor{dark-blue}{rgb}{0.0,0.0,0.5}
\definecolor{RED}{rgb}{1.0,0.0,0.0}

% Some new command
\newcommand{\YAWTB}{\textsc{YAWTb}}
\newcommand{\YAWTBHomePage}{%
http://www.fyma.ucl.ac.be/projects/yawtb}
\newcommand{\Matlab}{%
\htmladdnormallink{Matlab}{http://www.mathworks.com}}

%
% Commands related to the authors
%
\newcommand{\ACoron}{%
\htmladdnormallink{Alain Coron}{mailto:alain.coron@free.fr}}
\newcommand{\LDemanet}{%
\htmladdnormallink{Laurent Demanet}{mailto:demanet@student.fsa.ucl.ac.be}}
\newcommand{\LJacques}{%
\htmladdnormallink{Laurent Jacques}{mailto:ljacques@fyma.ucl.ac.be}}
\newcommand{\ARivoldini}{%
\htmladdnormallink{Attilio Rivoldini}{mailto:rivoldini@fyma.ucl.ac.be}}
\newcommand{\PVandergheynst}{%
\htmladdnormallink{Pierre Vandergheynst}{mailto:Pierre.Vandergheynst@epfl.ch}}

% Mailing lists
\newcommand{\yawtbdevel}{\htmladdnormallink{yawtb-devel@lists.sourceforge.net}{%
mailto:yawtb-devel@lists.sourceforge.net}}
\newcommand{\yawtbcode}{\htmladdnormallink{yawtb-code@lists.sourceforge.net}{%
mailto:yawtb-code@lists.sourceforge.net}}

% Extra commands that one could use in the doc
\newcommand{\libfun}[1]{\texttt{#1}}  % \libfun : Library function
\newcommand{\libvar}[1]{\texttt{#1}}  % \libvar : Library variable


\begin{htmlonly}
  \pagecolor[named]{White}
\end{htmlonly}


\pagestyle{headings}
\setlength{\parindent}{0cm}

\begin{document}

\title{%
Yet Another Wavelet Toolbox (YAW Toolbox) Reference Guide\\ \medskip
  Version 0.1.1}
\author{%
  The YAW Toolbox Team}

\begin{titlepage}
\begin{center}
\includegraphics[width=2cm,keepaspectratio]{yawtb_cwtsph.eps}\quad
\includegraphics[width=11cm,keepaspectratio]{logo_bgwhite.eps}\\
\vspace{0.5cm}
{\LARGE\bf Yet Another Wavelet Toolbox\\[0.5cm]
Reference Guide}\\[1cm] 
Version 0.1.1\\[0.5cm]
The YAW Toolbox Team\\[0.5cm]
\today
\end{center}
\end{titlepage}

%% VERSO TITLEPAGE IN LATEX
\begin{latexonly}
\begin{titlepage}
\clearpage
\vspace*{17cm}
\thispagestyle{empty}
Copyright \copyright\ 2001-2002,  The YAW Toolbox Team.

\medskip

Permission is granted to copy, distribute and/or modify this document under
the terms of the GNU Free Documentation License, Version 1.1 or any later
version published by the Free Software Foundation; with the Invariant
Sections being ``GNU General Public License'', with no Front-Cover Texts
and with no Back-Cover Texts. A copy of the license is included in the
section entitled ``GNU Free Documentation License''.


% (a) The Front-Cover text is: ``Yet Another Wavelet Toolbox (YAW Toolbox)
% Reference Guide''.
%
% (b) The Back-Cover text is: `` ``.
\end{titlepage}
\end{latexonly}

%% VERSO TITLEPAGE IN HTML 
\begin{htmlonly}
\ \\[2cm]
\htmlrule
\ \\[1cm]
Copyright \copyright\ 2001-2002,  The YAW Toolbox Team.

\medskip

Permission is granted to copy, distribute and/or modify this document under
the terms of the GNU Free Documentation License, Version 1.1 or any later
version published by the Free Software Foundation; with the Invariant
Sections being ``GNU General Public License'', with no Front-Cover Texts
and with no Back-Cover Texts. A copy of the license is included in the
section entitled ``GNU Free Documentation License''.


% (a) The Front-Cover text is: ``Yet Another Wavelet Toolbox (YAW Toolbox)
% Reference Guide''.
% 
% (b) The Back-Cover text is: `` ``.
\end{htmlonly}

\htmlrule

\abstract{The aim of this toolbox is to provide a clear and well documented
  implementation in Matlab of some time-frequency and time-scale
  transformations like the well known continuous/discrete wavelet
  transforms.
  
  In this guide each function of the Yet Another Wavelet Toolbox is
  documented.}

%\htmlrule

\tableofcontents


\section{History and contributors}
\label{sec:HisCon}

The \htmladdnormallink{\YAWTB{}}{\YAWTBHomePage}{} project started on
December 2000 few days after \LJacques' decisive E-Mail.  \ACoron,
\LJacques, \ARivoldini, and \PVandergheynst{} make up the core of the team.
Later \LDemanet{} joined the team.

% It is registered as follows:
%   + Project Full Name:  YAW Toolbox
%   + Project Unix Name:  yawtb
%   + CVS Server:         cvs.yawtb.sourceforge.net
%   + Shell/Web Server:   yawtb.sourceforge.net

The software is distributed under GNU General Public License.  So your
contributions are welcome.  You may join the two following mailing lists:
\begin{itemize}
\item \yawtbdevel{} is devoted to the development.  We discuss on features
  that we plan to add and how to code them.
\item \yawtbcode{} is devoted to the users.  They could
  send remarks, submit patches or new M-file.
\end{itemize}


\section{How to contribute ?}
\label{sec:HowToCon}

\subsection{Writing the Matlab headers}
\label{sec:WriMatHea}


\begin{itemize}
\item First of all, you should use the template that you will find in\\
  \texttt{tools/devel/xemacs\_macro/mfile\_header\_part.txt}.
\item A \verb|\YAWTB| command exist to denote the \YAWTB{} toolbox
\item Short description are better.  We do not need to say the one function
  compute or return something.  For example,
  \begin{itemize}
  \item cwt1d: 1D continuous wavelet transform
  \item gauss1d: 1D Gaussian wavelet in the frequency domain
  \end{itemize}
\item Some commands need and underscore \_.  In \LaTeX{} the \verb|_| has a
  special meaning.  Protect this character.
\item Use \verb|\libfun{function}| if you means the function
  \libfun{function}
\item Use \verb|\libvar{variable}| if you means the variable \libvar{variable}
\item In ths Sections ``Input Data'' and ``Output Data'' we indicate the
  type of the variable.  We already use the following types:
  \begin{itemize}
  \item BOOLEAN : A real number with two possible values $0$ or $1$
  \item COMPLEX MATRIX
  \item CPLX a complex
  \item CPLX MATRIX
  \item CPLX VECTOR a vector of complexes
  \item MATRIX : the matrix may be real or complex
  \item MISC when the type depends on the context
  \item REAL
  \item REAL SCALAR
  \item REAL MATRIX
  \item REAL VECTOR
  \item SCALAR : real or complex scalar
  \item STRING
  \item STRUCT a \Matlab structure
  \item VECTOR
  \item YAWTB OBJECT a structure that stores the output of functions like
    \libfun{cwt1d}
  \end{itemize}
  Some are redundant.  I guess that some should be considered as obsolete.
\item The environments 
\begin{verbatim}
     \begin{latexonly} ... \end{latexonly}
\end{verbatim} or
\begin{verbatim}
      \begin{htmlonly} ... \end{htmlonly}
\end{verbatim} 
   enable you to include different comments in the printed and htmlize
   versions.
 \item Each function will get its own label during processing the doc.  This
   label is build with the string \texttt{cmd:} followed by the name of the
   Matlab function.  For example the label of the \libfun{cauchy2d} command is
   \texttt{cmd:cauchy2d}.  Even the subsections of each command have their
   own labels.  So you can reference the functions with the command
   \texttt{\\htmlref}.  This command has two mandatory argument, the text
   that will appeak as a link and the reference of the link.
 \item Fill in the ``See Also'' sections with Perl regular expression
   separated by spaces.  The regular expression should match the name of
   files (without the Matlab extension) in the tree of the toolbox on a Unix
   machine.  Those regular will be expanded accordingly by a Perl script.
   
   For example, if you wish to add the function cwt2d, you may add
   \texttt{$\hat{ }$cwt2d\$} or \texttt{$\hat{ }$continuous/.*/cwt2d\$}.  If
   you wish to add all the functions under the continuous directory, type
   \texttt{$\hat{ }$continuous/.*}.

   If the regular expression does not match any name in the \YAWTB{} tree, then
   the leading $\hat{}$ (if any) and the terminated \$ (if any) of the expression
   are removed.  The remaining expression is added to the list of files.  So
   you may add the \texttt{fft} command by typing \texttt{$\hat{}$fft\$}.

\end{itemize}

\subsection{TODO}
\label{sec:TODO}

I list a lot of possible improvements:
\begin{itemize}
\item Parse the Matlab header files and build a tree to be able to perform
  easy transformation.  Having a tree should help at applying different type
  of transformation (depending on the context) on the different section
  (Example section, See Also section).  But building a tree may require a
  lot of memory.
\item Fill in the \htmlref{Section ``Writing the Matlab headers''}{sec:WriMatHea}
% \item The ``See Also'' should accept regex as command and replace those
%  regex by a direct link to the documentation.  It also means that we should
%  automatically add a label to each command.
% \item Insert an index.  In this index, the index entry Wavelet 1D could
%   point to all the 1D implemented wavelets
\item Code example Section: We should be able to extract the example code
  from the documentation, run separately those pieces of code and include
  the textual output as well as the figures in the final documentation.
  There should be a line-numbering system of the code and an automatic
  labelling of the code and the figure.

  How could we implement this?
  \begin{itemize}
  \item Maybe with a code environement like this:
\begin{verbatim}
  \begin{code}[number=<number_of_figures>,
               label=<label>]
  \end{code}
\end{verbatim}
  \item I suggest to internally construct the label of the code with
    \texttt{cmd:$<$cmd\_name$>$:code:$<$label$>$}
  \item The label of the associated figures could be\\
    \texttt{cmd:$<$cmd\_name$>$:code:$<$label$>$:fig$<$fig\_number$>$}
  \item The Matlab command line input/output\\
    \texttt{$<$cmd\_name$>$:code:$<$label$>$:Matlab.txt}.  Which directory?
  \item The figures could be saved under the name\\
    \texttt{$<$cmd\_name$>$:code:$<$label$>$:fig$<$fig\_number$>$.$<$extension$>$}.
    Which directory?
  \end{itemize}

\end{itemize}

\clearpage

%
% The commands sorted by category
%
\section{The commands by category}
\label{sec:ComCat}

\subsection{Contiuous Wavelet Transform}

\subsubsection{1D}


\input{TMP_cmd_cont_1d}



\subsubsection{1D+T}

\input{TMP_cmd_cont_1dt}
 
\subsubsection{2D}

\input{TMP_cmd_cont_2d}

\subsubsection{3D}

\input{TMP_cmd_cont_3d}

\subsubsection{Sphere}

\input{TMP_cmd_cont_sph}


\subsection{Discrete Wavelet Transform}

% \subsubsection{Frames 1D}

% \input{TMP_cmd_disc_frames_1d}


\subsubsection{Frames 2D}

\input{TMP_cmd_disc_frames_2d}

\subsubsection{Spherical Frames}

\input{TMP_cmd_disc_frames_sph}

% subsubsection{Laplacian}

% \input{TMP_cmd_disc_lap}


% \subsubsection{Matching Pursuit}

% \input{TMP_cmd_disc_mat}


% \subsubsection{Orthogonal}

% \input{TMP_cmd_disc_ort}

\subsubsection{Packet}

\paragraph{2d}

\input{TMP_cmd_disc_pack_2d}



\subsection{Interfaces}

\subsubsection{SpharmonicKit}

\input{TMP_cmd_interfaces_spharmonickit}



\subsection{Sample}

\input{TMP_cmd_samp}


\subsection{Help, help !}

\input{TMP_cmd_help}

\subsection{Tools}

\input{TMP_cmd_tool}



\clearpage

\section{The commands in alphabetical order}
\label{sec:ComAlpOrd}

\input{TMP_cmd_by_alpha}

% Include the GNU GENERAL PUBLIC LICENSE
\input{GPL}

% Include the GNU Free Documentation License
\input{FDL}

\end{document}
